% ============================================================================
%  Capitolo 5 -- Rappresentazione dei numeri: gli interi         (Lezione 05)
% ============================================================================
\chapter{Rappresentazione dei numeri: gli interi}
\label{ch:integers}

\epigraphline{Conterò poco, è vero: --- diceva l'Uno ar Zero --- ma tu che
vali? Gnente: propio gnente.}{Trilussa, \emph{Nummeri}, 1944}

\begin{chapgoals}
Alla fine di questo capitolo saprete
\begin{itemize}
\item convertire interi e frazioni tra le basi 2, 8, 10 e 16;
\item dire quali numeri naturali stanno in $n$ bit e riconoscere un overflow;
\item codificare e decodificare interi con segno in segno e modulo e in complemento a due;
\item sommare, sottrarre e cambiare di segno in complemento a due, e rilevare l'overflow;
\item spiegare perché $0-1$ appare come \hex{FFFFFFFF} in un registro \riscv;
\item scegliere l'istruzione \riscv giusta per dati con segno e senza segno.
\end{itemize}
\end{chapgoals}

% ----------------------------------------------------------------------------
\section{Riscaldamento sulla macchina}

Un marcatore sta sullo schermo, alla posizione 2 di una linea orizzontale. Il
programma legge il joystick e, finché \key{LEFT} è premuto, sottrae uno dalla
posizione a ogni frame (Listato~\ref{lst:int:warm}); la funzione \code{draw}
stampa la posizione in esadecimale con \code{prg32\_console\_hex32}.

\asmfile[caption={Tenete premuto \key{LEFT} e guardate il numero (adattato dal laboratorio
su stato e ingresso del libro di \prg).},label=lst:int:warm]{code/l05/c03.S}

\begin{prgtry}
Tenete premuto \key{LEFT}. Il valore mostrato fa \texttt{00000002},
\texttt{00000001}, \texttt{00000000} \ldots\ e poi \texttt{FFFFFFFF}, e il
marcatore scompare. (Sulla scheda e in QEMU la console è disegnata sullo
schermo; \prgqt la raccoglie nel suo log.) È un errore? \hex{FFFFFFFF} è il
numero 4\,294\,967\,295 oppure il numero $-1$?
\end{prgtry}

La risposta è: entrambi. Dipende da come si leggono i bit, e questo capitolo dà
le regole.

% ----------------------------------------------------------------------------
\section{Basi e conversioni}

\subsection{Sistemi posizionali}

\begin{definition}{Rappresentazione posizionale}
Sia $b\ge 2$ un intero, la \term{base}. Una stringa di cifre
$d_{n-1}\,d_{n-2}\cdots d_1\,d_0$, con $0\le d_i\le b-1$ per ogni $i$,
rappresenta in base $b$ il numero naturale
\begin{equation}\label{eq:int:positional}
N \;=\; \sum_{i=0}^{n-1} d_i\, b^{\,i}
  \;=\; d_{n-1}b^{\,n-1}+\dots+d_1 b+d_0 .
\end{equation}
La cifra $d_{n-1}$ è la \term{più significativa}, $d_0$ la \term{meno
significativa}.
\end{definition}

In informatica si usano quattro basi (Tabella~\ref{tab:int:bases}). Ogni numero
naturale ha esattamente una rappresentazione in ciascuna base, a parte gli zeri
aggiunti a sinistra.

\begin{table}[htb]
\centering\small
\caption{Le quattro basi dell'informatica.}\label{tab:int:bases}
\begin{tabular}{@{}clll@{}}
\toprule
\textbf{Base $b$} & \textbf{Cifre} & \textbf{Nome} & \textbf{Tredici si scrive}\\
\midrule
2 & 0, 1 & binaria & \texttt{1101}\\
8 & 0, 1, \ldots, 7 & ottale & \texttt{15}\\
10 & 0, 1, \ldots, 9 & decimale & \texttt{13}\\
16 & 0, \ldots, 9, A, \ldots, F & esadecimale & \texttt{D}\\
\bottomrule
\end{tabular}
\end{table}

\subsection{Le basi che sono potenze di due condividono gli stessi bit}

Se ogni cifra è scritta in binario puro, i bit di un numero sono gli stessi in
ogni base $b=2^{k}$: basta raggrupparli da destra, $k$ alla volta, completando
con zeri a sinistra. Per il numero 127:
\[
\texttt{01111111}_2 \;=\; \underbrace{\texttt{0111}}_{7}\,\underbrace{\texttt{1111}}_{F}
 \;=\; \hex{7F}
\qquad\text{e}\qquad
\texttt{01111111}_2 \;=\; \underbrace{\texttt{01}}_{1}\,\underbrace{\texttt{111}}_{7}\,\underbrace{\texttt{111}}_{7}
 \;=\; 177_8 .
\]
La ragione è che nella (\ref{eq:int:positional}) con $b=2$ i termini si possono
raccogliere $k$ alla volta, e ogni gruppo è una cifra moltiplicata per una
potenza di $2^{k}$.

\subsection{Convertire un intero: divisioni successive}

Per scrivere un numero naturale $N$ in base $r$: si divide $N$ per $r$; il
resto è la cifra successiva, a partire dalla meno significativa; si ripete sul
quoziente finché non è zero. I resti si leggono poi dall'ultimo al primo
(Tabella~\ref{tab:int:division}).

\begin{table}[htb]
\centering\small
\caption{Il numero 127 convertito per divisioni successive.}\label{tab:int:division}
\begin{tabular}{@{}rrr@{}}
\multicolumn{3}{c}{$r=8$}\\\toprule
$N$ & $N/8$ & $N \bmod 8$\\\midrule
127 & 15 & 7\\ 15 & 1 & 7\\ 1 & 0 & 1\\\bottomrule
\multicolumn{3}{c}{$127 = 177_8$}
\end{tabular}\qquad
\begin{tabular}{@{}rrr@{}}
\multicolumn{3}{c}{$r=16$}\\\toprule
$N$ & $N/16$ & $N \bmod 16$\\\midrule
127 & 7 & 15 = F\\ 7 & 0 & 7\\\bottomrule
\multicolumn{3}{c}{$127 = \hex{7F}$}
\end{tabular}\qquad
\begin{tabular}{@{}rrr@{}}
\multicolumn{3}{c}{$r=2$}\\\toprule
$N$ & $N/2$ & $N \bmod 2$\\\midrule
127 & 63 & 1\\ 63 & 31 & 1\\ 31 & 15 & 1\\ 15 & 7 & 1\\ 7 & 3 & 1\\ 3 & 1 & 1\\ 1 & 0 & 1\\\bottomrule
\multicolumn{3}{c}{$127 = 1111111_2$}
\end{tabular}
\end{table}

Il metodo è corretto perché, per la (\ref{eq:int:positional}),
$N = r\,(d_{n-1}r^{\,n-2}+\dots+d_1) + d_0$: il resto della divisione di $N$
per $r$ è $d_0$, e il quoziente è il numero le cui cifre sono
$d_{n-1}\cdots d_1$, a cui si applica lo stesso ragionamento. In C le due
operazioni sono \code{/} e \code{\%}; in \riscv, con l'estensione M, sono
\ins{divu} e \ins{remu}, e il Capitolo~\ref{ch:verm1p4} usa esattamente questo
algoritmo per stampare un punteggio.

\subsection{Convertire una frazione: moltiplicazioni successive}

Un numero $0\le N<1$ si scrive in base $r$ come
$N = d_{-1}r^{-1}+d_{-2}r^{-2}+\cdots$. Moltiplicando per $r$ si ottiene
$d_{-1} + (d_{-2}r^{-1}+\cdots)$: la parte intera del prodotto è la cifra
successiva, a partire dalla più significativa, e il procedimento si ripete
sulla parte frazionaria (Tabella~\ref{tab:int:fraction}).

\begin{table}[htb]
\centering\small
\caption{Due frazioni convertite in base 2.}\label{tab:int:fraction}
\begin{tabular}{@{}lll@{}}
\multicolumn{3}{c}{0,875: esatta}\\\toprule
Prodotto & Frazione & Cifra\\\midrule
$0,875\times 2$ & 0,75 & 1\\ $0,75\times 2$ & 0,5 & 1\\ $0,5\times 2$ & 0 & 1\\\bottomrule
\multicolumn{3}{c}{$0,875 = 0,111_2$}
\end{tabular}\qquad\qquad
\begin{tabular}{@{}lll@{}}
\multicolumn{3}{c}{0,8: periodica}\\\toprule
Prodotto & Frazione & Cifra\\\midrule
$0,8\times 2$ & 0,6 & 1\\ $0,6\times 2$ & 0,2 & 1\\ $0,2\times 2$ & 0,4 & 0\\ $0,4\times 2$ & 0,8 & 0\\ $0,8\times 2$ & 0,6 & 1 \ldots\\\bottomrule
\multicolumn{3}{c}{$0,8 = 0,\overline{1100}_2$}
\end{tabular}
\end{table}

\begin{prgwarn}
Una frazione finita in base 10 può essere periodica in base 2. Con un numero
finito di bit va troncata: da qui vengono gli errori di approssimazione dei
numeri reali (Capitolo~\ref{ch:reals}).
\end{prgwarn}

% ----------------------------------------------------------------------------
\section{Numeri naturali e overflow}

\subsection{Strategie di codifica}

\begin{definition}{Rappresentazione}
Una \term{rappresentazione} è una legge $X = R(x)$ che associa a un numero $x$
una stringa di bit $X$. Qui $x$ è il numero che si intende; $X$ è la
configurazione memorizzata nel registro.
\end{definition}

Il numero di bit dipende da quanti numeri si vogliono rappresentare; di solito
è un multiplo di 8, tipicamente 16, 32 o 64. La legge $R$ si sceglie secondo il
tipo di numero, naturale, con segno o razionale, e una buona legge è quella che
rende semplici da costruire come circuiti le operazioni aritmetiche e i
confronti.

\subsection{Numeri naturali su $n$ bit}

Con $n$ bit ci sono $2^{n}$ configurazioni, e la scelta naturale è associare a
ogni numero naturale da 0 a $2^{n}-1$ la stringa di bit che lo rappresenta nel
sistema posizionale binario. Con $n=4$ l'intervallo è $[0,15]$. Le dimensioni
trattate da \riscv sono nella Tabella~\ref{tab:int:naturals}.

\begin{table}[htb]
\centering\small
\caption{Numeri naturali su 8, 16 e 32 bit.}\label{tab:int:naturals}
\begin{tabular}{@{}rlll@{}}
\toprule
\textbf{Bit} & \textbf{Intervallo} & \textbf{In \riscv} & \textbf{Massimo}\\
\midrule
8 & $0\ldots 2^{8}-1$ & byte (\ins{lbu}) & 255\\
16 & $0\ldots 2^{16}-1$ & mezza parola (\ins{lhu}) & 65\,535\\
32 & $0\ldots 2^{32}-1$ & parola, registro & 4\,294\,967\,295\\
\bottomrule
\end{tabular}
\end{table}

Il numero 4\,294\,967\,295 è la lettura di \hex{FFFFFFFF} come numero naturale:
la prima metà della risposta al riscaldamento.

\subsection{Proprietà delle stringhe binarie}

Dalla (\ref{eq:int:positional}) seguono subito tre proprietà.

\begin{enumerate}
\item \emph{L'estensione a sinistra con zeri non cambia il valore}:
  $5 = 101_2 = 00101_2$.
\item \emph{Moltiplicare per $2^{k}$ è uno scorrimento a sinistra di $k$
  posizioni}: $00101_2\times 2 = 01010_2$, cinque diventa dieci. Ogni cifra
  $d_i$ passa alla posizione $i+k$, dove pesa $2^{k}$ volte tanto.
\item \emph{La divisione intera per $2^{k}$ è uno scorrimento a destra di $k$
  posizioni}: $00101_2/2 = 00010_2$, cinque diventa due, e il resto va perso.
\end{enumerate}

Gli scorrimenti costano molto meno delle moltiplicazioni, e i compilatori li
usano ogni volta che possono (Listato~\ref{lst:int:shifts}). L'istruzione
\code{slli t1, t1, 2} del Capitolo~\ref{ch:information} moltiplicava così un
indice per quattro. Gli scorrimenti sono anche la base dell'aritmetica in
virgola fissa (Capitolo~\ref{ch:reals}).

\asmfile[caption={Gli scorrimenti come moltiplicazioni e divisioni.},label=lst:int:shifts]{code/l05/c04.S}

\subsection{Overflow}

\begin{definition}{Overflow}
Si ha un \term{overflow} quando il risultato di un'operazione non è contenuto
nel sottoinsieme finito di numeri che i bit disponibili possono rappresentare.
\end{definition}

I registri hanno una dimensione finita, e ce n'è un numero finito. Gli insiemi
numerici sono infiniti, e di ciascuno si può rappresentare solo un sottoinsieme
finito. Gli operatori aritmetici sono chiusi sull'intero insieme, non sul
sottoinsieme finito. Con 7 bit si codifica l'intervallo $[0,127]$: 100 e 38
sono rappresentabili, la loro somma 138 no.

\subsection{Aritmetica su quattro bit}

L'aritmetica binaria segue le stesse regole colonna per colonna
dell'aritmetica decimale, con l'unica tabellina $0+0=0$, $0+1=1$, $1+1=0$ con
riporto di~1.

\begin{center}\ttfamily
\begin{tabular}{@{}r@{\quad}r@{}}
\multicolumn{2}{c}{\textsf{\small Somma}}\\
6 + & 0110 +\\ 8 = & 1000 =\\\hline 14 & 1110
\end{tabular}\qquad
\begin{tabular}{@{}r@{\quad}r@{}}
\multicolumn{2}{c}{\textsf{\small Somma con overflow}}\\
14 + & 1110 +\\ 3 = & 0011 =\\\hline 17 & \textcolor{prgred}{1}\,0001
\end{tabular}\qquad
\begin{tabular}{@{}r@{\quad}r@{}}
\multicolumn{2}{c}{\textsf{\small Differenza}}\\
11 - & 1011 -\\ 5 = & 0101 =\\\hline 6 & 0110
\end{tabular}\qquad
\begin{tabular}{@{}r@{\quad}r@{}}
\multicolumn{2}{c}{\textsf{\small Prodotto}}\\
5 x & 0101 x\\ 3 = & 0011 =\\\hline 15 & 000\,1111
\end{tabular}
\end{center}

Nella seconda somma serve un quinto bit: il registro di quattro bit conserva
\texttt{0001} e il riporto va perso. I processori segnalano un riporto perso in
modi diversi: molti hanno un \emph{flag di riporto} (carry flag); \riscv non ne
ha e lascia che sia il programma a calcolarlo con un confronto, come mostra la
Sezione~\ref{sec:int:back}. Il prodotto di due numeri di $n$ bit richiede fino
a $2n$ bit: su RV32 l'istruzione \ins{mul} dà i 32 bit bassi e \ins{mulhu} i 32
alti.

% ----------------------------------------------------------------------------
\section{Interi con segno}

Ci sono tre modi di scrivere un segno meno con i bit. Due sono presentati qui
e il terzo, con una variante, nel Capitolo~\ref{ch:reals}.

\subsection{Segno e modulo}

Il bit più significativo è il segno (1 per i negativi) e gli altri $n-1$ bit
sono il modulo:
\[
X = |x| + 2^{\,n-1}s,\qquad s=\begin{cases}1 & x<0\\ 0 & x\ge 0.\end{cases}
\]
L'intervallo va da $-(2^{n-1}-1)$ a $+(2^{n-1}-1)$. Lo zero ha due
rappresentazioni, $+0$ e $-0$, quindi si rappresentano solo $2^{n}-1$ numeri. E
l'aritmetica deve lavorare separatamente su segno e modulo: per sommare due
numeri di segno opposto bisogna trovare il modulo maggiore, sottrarne il
minore, e dare al risultato il segno dell'operando con il modulo maggiore. La
rappresentazione è semplice per le persone e scomoda per i circuiti.
Sopravvive nella mantissa dei numeri in virgola mobile
(Capitolo~\ref{ch:reals}).

\subsection{Complemento a due}

\begin{definition}{Complemento a due}
Su $n$ bit, la rappresentazione in \term{complemento a due} di un intero $x$
con $-2^{\,n-1}\le x\le 2^{\,n-1}-1$ è il numero naturale
\[
X=\begin{cases} x & \text{se } 0\le x<2^{\,n-1},\\[2pt]
2^{\,n}-|x| & \text{se } -2^{\,n-1}\le x<0,\end{cases}
\]
scritto in binario su $n$ bit. In modo equivalente, $X = x \bmod 2^{\,n}$.
\end{definition}

L'intervallo va da $-2^{n-1}$ a $+2^{n-1}-1$, e lo zero ha una sola
rappresentazione. Il bit più significativo vale 1 se e solo se il numero è
negativo, perché i numeri negativi occupano le configurazioni da $2^{n-1}$ a
$2^{n}-1$. È la rappresentazione usata da ogni processore moderno, \riscv
compreso: un registro di 32 bit letto come intero con segno copre
$-2^{31}\ldots 2^{31}-1$, da circa meno due miliardi a circa più due miliardi.
L'intervallo è asimmetrico, con un numero negativo in più rispetto ai positivi,
e questo ha una conseguenza che comparirà tra poco.

\subsection{Quattro bit, tre letture}

La Tabella~\ref{tab:int:readings} legge ciascuna delle sedici configurazioni di
quattro bit nei tre modi visti finora. Quando il bit più significativo vale 0
le tre letture coincidono; quando vale 1 differiscono.

\begin{table}[htb]
\centering\small
\caption{Quattro bit, tre letture. S-M: segno e modulo; C2: complemento a due.}\label{tab:int:readings}
\begin{tabular}{@{}cccc@{\qquad}cccc@{}}
\toprule
\textbf{Bit} & \textbf{Senza segno} & \textbf{S-M} & \textbf{C2} &
\textbf{Bit} & \textbf{Senza segno} & \textbf{S-M} & \textbf{C2}\\
\midrule
\texttt{0000} & 0 & $+0$ & 0 & \texttt{1000} & 8 & $-0$ & $-8$\\
\texttt{0001} & 1 & $+1$ & $+1$ & \texttt{1001} & 9 & $-1$ & $-7$\\
\texttt{0010} & 2 & $+2$ & $+2$ & \texttt{1010} & 10 & $-2$ & $-6$\\
\texttt{0011} & 3 & $+3$ & $+3$ & \texttt{1011} & 11 & $-3$ & $-5$\\
\texttt{0100} & 4 & $+4$ & $+4$ & \texttt{1100} & 12 & $-4$ & $-4$\\
\texttt{0101} & 5 & $+5$ & $+5$ & \texttt{1101} & 13 & $-5$ & $-3$\\
\texttt{0110} & 6 & $+6$ & $+6$ & \texttt{1110} & 14 & $-6$ & $-2$\\
\texttt{0111} & 7 & $+7$ & $+7$ & \texttt{1111} & 15 & $-7$ & $-1$\\
\bottomrule
\end{tabular}
\end{table}

\begin{prgconcept}
I bit non dicono quale sia la colonna giusta. Lo dice il programma, scegliendo
le sue istruzioni.
\end{prgconcept}

\subsection{La ruota dei numeri}

Con $n$ bit l'aritmetica è \emph{modulo $2^{n}$}: la retta dei numeri si chiude
in un cerchio (Figura~\ref{fig:int:wheel}). Sommare 1 sposta di un passo in
senso orario, qualunque sia la lettura. Le due letture differiscono solo nel
punto in cui la ruota «si spezza». Nella lettura senza segno si spezza tra
\texttt{1111} e \texttt{0000}, e attraversare quel punto è un riporto:
$15+1=0$. Nella lettura in complemento a due si spezza tra \texttt{0111} e
\texttt{1000}, e attraversare quel punto è un overflow con segno: $7+1=-8$. Un
passo indietro da \texttt{0000} è \texttt{1111}: ecco perché $0-1$ dà tutti
uni.

\begin{figure}[htb]
\centering
\begin{tikzpicture}
  \draw[prgrule,line width=7mm] (0,0) circle (2.6);
  \foreach \k in {0,...,15}{
    \pgfmathsetmacro{\a}{90-\k*22.5}
    \pgfmathtruncatemacro{\s}{ifthenelse(\k>7,\k-16,\k)}
    \pgfmathtruncatemacro{\ba}{mod(int(\k/8),2)}
    \pgfmathtruncatemacro{\bb}{mod(int(\k/4),2)}
    \pgfmathtruncatemacro{\bc}{mod(int(\k/2),2)}
    \pgfmathtruncatemacro{\bd}{mod(\k,2)}
    \node[font=\ttfamily\scriptsize] at (\a:2.6) {\ba\bb\bc\bd};
    \node[font=\sffamily\small,text=prgsteel] at (\a:1.85) {\k};
    \ifnum\k>7
      \node[circle,fill=prgamber!30,inner sep=1pt,minimum size=7mm,font=\sffamily\small] at (\a:3.55) {$\s$};
    \else
      \node[circle,fill=prgaccent!20,inner sep=1pt,minimum size=7mm,font=\sffamily\small] at (\a:3.55) {$\s$};
    \fi
  }
  \draw[prgviolet,very thick,dashed] (101.25:2.1) -- (101.25:4.1) node[above left,lbl,text=prgviolet,align=right]{rottura senza segno:\\riporto};
  \draw[prgred,very thick,dashed] (-78.75:2.1) -- (-78.75:4.1) node[below right,lbl,text=prgred,align=left]{rottura con segno:\\overflow};
  \draw[->,thick,prgink] (80:1.1) arc (80:10:1.1) node[midway,below left,lbl]{$+1$};
\end{tikzpicture}
\caption{La ruota dei numeri per quattro bit. Numeri interni: valore senza segno; anello:
bit; dischi esterni: valore in complemento a due.}
\label{fig:int:wheel}
\end{figure}

Lo stesso sommatore serve a entrambe le letture; cambia solo il punto in cui la
ruota si spezza.

\subsection{Aritmetica in complemento a due}

\begin{theorem}{Addizione in complemento a due}
Siano $x$, $y$ e $x+y$ tutti in $[-2^{n-1},\,2^{n-1}-1]$, e siano $X$, $Y$ le
rappresentazioni di $x$ e $y$. Allora la rappresentazione di $x+y$ è
$(X+Y)\bmod 2^{n}$: la somma delle due configurazioni come numeri naturali,
scartando il riporto uscente dal bit più significativo.
\end{theorem}

\noindent\emph{Dimostrazione.} $X\equiv x$ e $Y\equiv y \pmod{2^{n}}$ per
definizione, quindi $X+Y\equiv x+y \pmod{2^{n}}$; e la rappresentazione di
$x+y$ è, per definizione, $(x+y)\bmod 2^{n}$. \hfill$\square$

\smallskip
L'addizione non ha quindi bisogno di conoscere i segni: il sommatore per i
numeri naturali dà già la risposta giusta.

\begin{center}\ttfamily
\begin{tabular}{@{}r@{\quad}r@{}}
\multicolumn{2}{c}{\textsf{\small Segni opposti}}\\
2 + & 0010 +\\ -6 = & 1010 =\\\hline -4 & 1100
\end{tabular}\qquad\qquad
\begin{tabular}{@{}r@{\quad}r@{}}
\multicolumn{2}{c}{\textsf{\small Due negativi}}\\
-2 + & 1110 +\\ -3 = & 1101 =\\\hline -5 & \textcolor{prgcomment}{1}\,1011
\end{tabular}\qquad\qquad
\begin{tabular}{@{}r@{\quad}r@{}}
\multicolumn{2}{c}{\textsf{\small Prodotto}}\\
2 x & 0010 x\\ -3 = & 1101 =\\\hline -6 & \textcolor{prgcomment}{001}\,1010
\end{tabular}
\end{center}

Lo stesso ragionamento vale per la moltiplicazione: $XY\equiv xy\pmod{2^{n}}$,
quindi gli $n$ bit bassi di un prodotto sono corretti qualunque siano i segni.
Nel terzo esempio $2\times 13 = 26 = 1\,1010_2$, e i quattro bit bassi
\texttt{1010} si leggono $-6$. Per questo \riscv ha una sola \ins{mul} per
operandi con segno e senza segno.

\subsection{Cambiare di segno un numero}

Per definizione $-x$ è rappresentato da $2^{n}-X$. Due metodi lo calcolano
senza una sottrazione.

\begin{description}[style=nextline,leftmargin=1.5em]
\item[Metodo 1: complementare ogni bit, poi sommare 1.]
  $6=\texttt{0110}$, complemento \texttt{1001}, più uno: $\texttt{1010}=-6$.
  Funziona perché il complemento di ogni bit di $X$ è il numero $(2^{n}-1)-X$,
  una sottrazione da tutti uni in cui non può esserci alcun prestito; sommando
  1 si ottiene $2^{n}-X$.
\item[Metodo 2: scorrere da destra.] Partendo da destra, si copiano i bit fino
  al primo 1 compreso, e si complementano tutti gli altri. $5=\texttt{0101}$ dà
  $\texttt{1011}=-5$; $1=\texttt{0001}$ dà $\texttt{1111}=-1$. È il metodo~1
  con l'addizione svolta a mente: il $+1$ finale si propaga attraverso gli uni
  finali del complemento, che sono gli zeri finali dell'originale.
\end{description}

Cambiare di segno due volte restituisce il numero originale, con un'eccezione:
il numero più negativo, $-2^{n-1}$, non ha un compagno positivo, e il suo
cambio di segno va in overflow e restituisce il numero stesso.

\asmfile[caption={Il cambio di segno in \riscv.},label=lst:int:neg]{code/l05/c05.S}

\begin{prgconcept}[La sottrazione è un'addizione]
$y - x = y + (2^{n}-x) \pmod{2^{n}}$. Un sommatore e un complementatore sono
gli unici componenti necessari per costruire addizione, sottrazione,
moltiplicazione e divisione.
\end{prgconcept}

È la ragione ingegneristica del complemento a due: un solo circuito, il
sommatore del Capitolo~\ref{ch:combinational}, fa tutto il lavoro. Nel datapath
di \riscv \ins{sub} è \ins{add} con il secondo operando complementato e un
riporto entrante di~1.

\subsection{Rilevare l'overflow}

\begin{center}\ttfamily
\begin{tabular}{@{}r@{\quad}r@{}}
\multicolumn{2}{c}{\textsf{\small Senza segno, 8 bit}}\\
128 + & 1000 0000 +\\ 130 = & 1000 0010 =\\\hline 258 & \textcolor{prgred}{1} 0000 0010
\end{tabular}\qquad\qquad\qquad
\begin{tabular}{@{}r@{\quad}r@{}}
\multicolumn{2}{c}{\textsf{\small Con segno, 8 bit}}\\
127 + & 0111 1111 +\\ 127 = & 0111 1111 =\\\hline 254 & \textcolor{prgred}{1}111 1110
\end{tabular}
\end{center}

Nella prima somma c'è un riporto uscente, e il registro contiene 2, la somma
modulo 256. Nella seconda non c'è riporto uscente, ma il risultato si legge
$-2$.

\begin{theorem}{Condizioni di overflow}
In un'addizione su $n$ bit:
\begin{enumerate}[nosep]
\item si ha un overflow \emph{senza segno} se e solo se c'è un riporto uscente
  dal bit più significativo;
\item si ha un overflow \emph{con segno} se e solo se due operandi con lo
  stesso segno danno una somma di segno opposto: due positivi danno un
  negativo, o due negativi danno un positivo.
\end{enumerate}
La somma di due numeri di segno opposto non va mai in overflow.
\end{theorem}

L'ultima affermazione vale perché una tale somma è compresa tra i due operandi,
entrambi rappresentabili. Le due condizioni sono indipendenti: la stessa
addizione può far scattare l'una, entrambe o nessuna. In termini di circuiti la
condizione con segno equivale a questa: il riporto \emph{entrante} nel bit più
significativo è diverso dal riporto \emph{uscente} da esso.

% ----------------------------------------------------------------------------
\section{Ritorno alla macchina}\label{sec:int:back}

\subsection{Il riscaldamento, spiegato}

La variabile \code{marker\_x} conteneva 0. L'istruzione
\code{addi t2, t2, -1} ha fatto un passo indietro su una ruota con $2^{32}$
posizioni, e il registro ora contiene trentadue uni: \hex{FFFFFFFF}. Letto
senza segno è 4\,294\,967\,295; letto in complemento a due è $-1$. Il sommatore
ha fatto la stessa cosa in entrambi i casi. La funzione
\code{prg32\_console\_hex32} stampa la configurazione, non un'interpretazione:
non ha né ragione né torto. E il marcatore scompare perché a
\code{prg32\_gfx\_rect} si chiede di disegnare a una $x$ enorme, oppure a
$-1$: fuori dallo schermo in ogni caso (Figura~\ref{fig:int:marker}). Non è un
errore del processore; è una decisione che manca nel programma.

\begin{figure}[htb]
\centering
\includegraphics[width=0.5\textwidth]{lab2_step2_move.png}
\caption{Il programma del laboratorio: sette lampade per i bit dei pulsanti e un
marcatore.}
\label{fig:int:marker}
\end{figure}

\subsection{Con segno o senza segno? Sceglie l'istruzione}

I registri di \riscv non hanno un tipo. Gli stessi 32 bit sono con segno o
senza segno solo agli occhi dell'istruzione che li legge
(Tabella~\ref{tab:int:signed}).

\begin{table}[htb]
\centering\small
\caption{Istruzioni con segno e senza segno.}\label{tab:int:signed}
\begin{tabularx}{\textwidth}{@{}lllL@{}}
\toprule
\textbf{Operazione} & \textbf{Con segno} & \textbf{Senza segno} & \textbf{Perché differiscono}\\
\midrule
somma, sottrazione & \ins{add}, \ins{addi}, \ins{sub} & le stesse & complemento a due: un sommatore per entrambi\\
prodotto (parola bassa) & \ins{mul} & la stessa & i 32 bit bassi non dipendono dal segno\\
confronto & \ins{slt}, \ins{slti} & \ins{sltu}, \ins{sltiu} & \hex{FFFFFFFF} è minore di 0? sì, se con segno\\
salto condizionato & \ins{blt}, \ins{bge} & \ins{bltu}, \ins{bgeu} & la stessa domanda, come salto\\
scorrimento a destra & \ins{sra}, \ins{srai} & \ins{srl}, \ins{srli} & copia il bit di segno, o inserisce zeri\\
caricamento di un byte & \ins{lb} & \ins{lbu} & estensione del segno o con zeri\\
divisione, resto & \ins{div}, \ins{rem} & \ins{divu}, \ins{remu} & i moduli si leggono in modo diverso\\
\bottomrule
\end{tabularx}
\end{table}

\subsection{Tenere il marcatore sullo schermo}

La decisione che mancava è presa nel Listato~\ref{lst:int:clamp}: dopo lo
spostamento, la posizione viene \term{limitata} (clamp) all'intervallo
$[0, 304]$.

\asmfile[caption={Limitare la posizione (laboratorio su stato e ingresso, passo 3).},label=lst:int:clamp]{code/l05/c08.S}

\ins{bgez} legge \reg{t2} come numero con segno: esamina il bit di segno. Anche
\ins{ble} è un confronto con segno: $-1$ è minore di 304, mentre con \ins{bleu}
la configurazione \hex{FFFFFFFF} sarebbe di gran lunga maggiore. La costante
\code{-SPEED} è un immediato: dodici bit in complemento a due, da $-2048$ a
$+2047$.

\begin{prgconcept}
Un gioco non rileva l'overflow a cose fatte: mantiene i suoi valori dentro un
intervallo in cui l'overflow non può avvenire.
\end{prgconcept}

\subsection{Rilevare l'overflow senza flag}

\riscv non ha un flag di riporto né un flag di overflow: meno stati nascosti,
un processore più semplice. Quando un programma ha davvero bisogno di saperlo,
lo chiede con un confronto (Listato~\ref{lst:int:ovf}). Per i numeri senza
segno, una somma che ha fatto il giro è sempre minore di ciascuno dei due
operandi. Per i numeri con segno, sommare un numero negativo deve rendere la
somma più piccola, e sommarne uno non negativo no. Sulla ruota, \ins{sltu}
rileva l'attraversamento della cima e il test con segno l'attraversamento del
fondo.

\asmfile[caption={I due test di overflow suggeriti dalla specifica \riscv.},label=lst:int:ovf]{code/l05/c09.S}

\subsection{Estensione del segno: lo stesso numero su più bit}

\begin{theorem}{Estensione del segno}
Un numero in complemento a due su $n$ bit si scrive su $m>n$ bit copiando il
suo bit di segno in tutte le nuove posizioni a sinistra.
\end{theorem}

\noindent\emph{Dimostrazione.} Se $x\ge 0$ il bit di segno è 0 e l'enunciato è
l'estensione con zeri dei numeri naturali. Se $x<0$ la sua configurazione su
$n$ bit è $X=2^{n}-|x|$ e la configurazione voluta è
$2^{m}-|x| = X + (2^{m}-2^{n})$; il numero $2^{m}-2^{n}$ è esattamente $m-n$
uni seguiti da $n$ zeri.
\hfill$\square$

\begin{center}\ttfamily
\begin{tabular}{@{}r@{\quad}l@{\quad}l@{}}
\textsf{$-1$ su 4 bit} & \phantom{1111\,}1111 & = -1\\
\textsf{$-1$ su 8 bit} & \textcolor{prgred}{1111}\,1111 & = -1\\
\textsf{$+5$ su 4 bit} & \phantom{0000\,}0101 & = +5\\
\textsf{$+5$ su 8 bit} & \textcolor{prgsteel}{0000}\,0101 & = +5\\
\end{tabular}
\end{center}

L'estensione con zeri trasformerebbe \texttt{1111} ($-1$) in
\texttt{0000\,1111} ($+15$). L'istruzione \code{addi t2, t2, -1} memorizza $-1$
in dodici bit, e il processore ne estende il segno a trentadue prima di
sommare. \ins{lb} estende il segno del byte che carica, \ins{lbu} riempie con
zeri: la scelta dipende da che cosa significa il byte.

% ----------------------------------------------------------------------------
% Generated by tools/textbook/export_listings.js from decks/plan.js and code/l05.S: do not edit.
\section{Il programma di questo capitolo}\label{sec:lab:l05}

I listati di questo capitolo sono parti di un solo programma,
\fname{code/l05.S} nel repository del corso, stampato per intero nel
Listato~\ref{lst:lab:l05}. I test automatici del materiale del corso lo
compilano e lo eseguono sull'esecutore senza interfaccia di \prgqt{}, con
ingressi programmati, e controllano il frame; la
Figura~\ref{fig:lab:l05} è il frame scritto dal test.

\begin{description}[style=nextline,leftmargin=1.5em]
\item[Sullo schermo.] I 32 bit di marker\_x come lampade. Parte da 2 (una lampada); ogni frame con LEFT premuto sottrae 1: dopo 2, 1 e 0 si accendono tutte le lampade, 0xFFFFFFFF, cioè -1.
\item[Che cosa premere.] Tenete premuto LEFT per qualche frame.
\end{description}

\begin{figure}[htb]
\centering
\includegraphics[width=0.5\textwidth]{lab_l05.png}
\caption{\fname{l05.S} su \prgqt{}.}\label{fig:lab:l05}
\end{figure}

\asmfile[caption={Il programma completo del capitolo, \fname{code/l05.S}.},label=lst:lab:l05]{code/l05/program.S}

Per eseguirlo su \prgqt{} i comandi sono quelli del
Listato~\ref{lst:lab:l05:run}; il Paragrafo~\ref{sec:prg:testing} dà gli
stessi passi per QEMU e per la scheda.

\begin{shell}[caption={Compilare \fname{l05.S} ed eseguirlo su \prgqt{}.},label=lst:lab:l05:run]
# 1. start PRG32-QT on this computer (it listens on port 8080)

# 2. in the PRG32 folder, with ESP-IDF active: build the cartridge
python3 -m prg32 cartridge build ~/arc12/code/l05.S --portable \
    --entry-prefix l05 --name l05 --out l05.prg32

# 3. send it to PRG32-QT and start it
python3 -m prg32 esp32c6 upload-and-run l05.prg32 \
    --url http://localhost:8080 --slot cart0

# 4. look, press, compare with the previous slide
\end{shell}


% ----------------------------------------------------------------------------
\section{Verificate di aver capito}

\begin{quiz}{Quale numero rappresenta in complemento a due la configurazione di 8 bit \texttt{1111 1100}?}
\item 252
\item $-4$
\item $-3$
\item $-124$
\end{quiz}

\begin{quiz}{Su 8 bit, quale addizione produce un overflow con segno?}
\item $100 + 27$
\item $100 + 28$
\item $-100 + 28$
\item $-1 + 1$
\end{quiz}

\begin{quiz}{\reg{t0} contiene \hex{FFFFFFFF} e \reg{t1} contiene 1. Dopo
\code{slt t2, t0, t1} e \code{sltu t3, t0, t1}:}
\item \reg{t2} = 1, \reg{t3} = 1
\item \reg{t2} = 0, \reg{t3} = 0
\item \reg{t2} = 1, \reg{t3} = 0
\item \reg{t2} = 0, \reg{t3} = 1
\end{quiz}

\begin{quiz}{Come si scrive 0,625 in binario?}
\item 0,101
\item 0,110
\item 0,625
\item 0,011
\end{quiz}

\section{Esercizi}

\begin{exercise}[prima su carta, poi su \prg]
\begin{enumerate}[nosep]
\item Convertite 2026 in binario, ottale ed esadecimale; convertite 0,1 in
  binario fino a otto bit frazionari.
\item Scrivete $-100$ su 8 bit in segno e modulo e in complemento a due.
\item Calcolate su 8 bit in complemento a due: $90+60$, $-90-60$, $90-60$.
  Quali vanno in overflow?
\item Nel codice del clamp sostituite \ins{bgez} e \ins{ble} con le loro
  versioni senza segno, prevedete il comportamento, poi eseguite.
\item Scrivete una \code{update} che somma \hex{10000000} a un contatore a ogni
  frame e accende un rettangolo rosso quando la somma con segno va in overflow.
\end{enumerate}
\expected{Su carta: $2026 = \hex{7EA} = 111\,1110\,1010_2$. Su \prg: il
rettangolo rosso lampeggia una volta ogni sedici frame, quando il contatore
passa da \hex{7\ldots} a \hex{8\ldots}.}
\end{exercise}

\section{Glossario}

\begin{glossarytable}
base (radix) & base\\
natural number & numero naturale\\
signed integer & numero relativo (intero con segno)\\
most significant bit & bit pi\`u significativo\\
remainder & resto\\
shift & scorrimento (shift)\\
two's complement & complemento a due (alla base)\\
sign and magnitude & segno e modulo\\
overflow & overflow (traboccamento)\\
carry / borrow & riporto / prestito\\
sign extension & estensione del segno\\
to clamp & limitare a un intervallo\\
\end{glossarytable}

\section{Conclusioni e letture}

Ogni base ha i suoi vantaggi: la base 2 per i circuiti, la base 16 per le
persone che leggono i bit. $n$ bit rappresentano $2^{n}$ numeri: quali, dipende
dalla legge di codifica. Il complemento a due trasforma la sottrazione in
addizione, così che un solo sommatore fa tutto. L'overflow è il prezzo dei
registri finiti: lo si rileva, oppure lo si evita per progetto. In \riscv i bit
non hanno un tipo: è l'istruzione a scegliere tra con segno e senza segno.

\paragraph{Letture.} \citet{patterson2020}, paragrafi 2.4 e 3.1--3.2;
\emph{Playing with RISC-V}, capitolo~6, «State and Input»
\citep{prg32textbook}; \citet{stallings2021}, i capitoli sui sistemi di
numerazione e sull'aritmetica del calcolatore.